% Part: many-valued-logic
% Chapter: syntax-and-semantics
% Section: formulas

\documentclass[../../../include/open-logic-section]{subfiles}

\begin{document}

\olfileid{mvl}{syn}{fml}

\olsection{\usetoken{P}{formula}}

\begin{defn}[فارمولا]
\ollabel{defn:formulas}
سیٹ~$\Frm[L]$ وچ قضیاتی زبان~$\Lang L$ دے \emph{!!{formula}s}
ہوندے نیں۔ ایہ سیٹ استقرا نال ایویں متعین ہوندا اے۔
\begin{enumerate}
\item ہر !!{propositional variable}~$\Obj p_i$ اک ایٹمی
  !!{formula} اے۔
\item ہر $0$-مقامی رابطہ، یعنی قضیاتی مستقل، جیہڑا~$\Lang L$
  وچ اے، اک ایٹمی !!{formula} اے۔
\item جے $\star$ اک $n$-مقامی رابطہ اے جیہڑا~$\Lang L$ وچ اے، تے $!A_1$،
\dots، $!A_n$، !!{formula}s نیں، تاں $\star(!A_1, \dots, !A_n)$
!!a{formula} اے۔
\tagitem{limitClause}{ہور کوئی شے !!a{formula} نہیں۔}{}
\end{enumerate}
جے $\star$ اک $1$-مقامی رابطہ اے، تاں $\star(!A_1)$ نوں
اکثر بس $\star !A_1$ لکھدے آں۔ جے $\star$ اک $2$-مقامی
رابطہ اے، تاں $\star(!A_1,!A_2)$ نوں اکثر $(!A_1 \star !A_2)$
لکھدے آں۔
\end{defn}

رواج مطابق اسیں اکثر سبھ توں باہر دیاں قوساں بِن دسے چھڈ
دیندے آں۔

\begin{ex}
  معیاری زبان~$\Lang{L_0}$ وچ $\Obj p_1 \lif (\Obj p_1
  \land \lnot \Obj p_2)$ اک فارمولا اے۔ حاصلِ ضرب والی منطق
  دی زبان وچ اینوں $\Obj p_1 \lif (\Obj p_1 \odot
  \lnot \Obj p_2)$ لکھیا جاندا اے۔ جے $1$-مقامی $\triangle$
  زبان وچ شامل کریے، تاں ساڈے کول $\triangle (\Obj p_1
  \land \Obj p_2) \lif (\triangle \Obj p_1 \land \triangle \Obj p_2)$
  ورگے فارمولے وی ہون گے۔
\end{ex}

\end{document}
